\documentclass[12pt]{article}

\usepackage{amsmath}
\usepackage{amssymb}
\usepackage{booktabs}
\usepackage{geometry}

\geometry{margin=1in}

\title{Lecture Scribe}
\date{}

\begin{document}

\maketitle
\textbf{Course:} CSE400 – Fundamentals of Probability in Computing\\
\textbf{Lecture:} L15–L18\\
\textbf{Topic:} Joint Random Variables, Joint Distributions, and Examples\\
\textbf{Instructor:} Dhaval Patel, PhD\\
\textbf{Institution:} SEAS, Ahmedabad University\\
\textbf{Date:} March 10, 2026

\hrule
\vspace{0.5cm}

\section*{1. Preliminaries}

\subsection*{Definitions Introduced in the Lecture}

\subsubsection*{Joint Random Variables}

When \textbf{two or more random variables are considered simultaneously}, they are called \textbf{joint random variables}.

Example representation:

\begin{itemize}
\item $X$ = first random variable
\item $Y$ = second random variable
\end{itemize}

These variables may be \textbf{correlated} or dependent on each other.

\subsection*{Notation and Symbols}

\begin{center}
\begin{tabular}{ll}
\toprule
Symbol & Meaning \\
\midrule
$X$ & First random variable \\
$Y$ & Second random variable \\
$p_{X,Y}(x,y)$ & Joint Probability Mass Function \\
$f_{X,Y}(x,y)$ & Joint Probability Density Function \\
$F_{X,Y}(x,y)$ & Joint Cumulative Distribution Function \\
\bottomrule
\end{tabular}
\end{center}

\hrule
\vspace{0.5cm}

\section*{2. Model Assumptions and Setup}

\subsection*{Motivation for Studying Joint Random Variables}

In many real-world systems, \textbf{multiple random variables occur simultaneously}, and their relationships must be analyzed together.

\subsection*{Examples Given in the Lecture}

\subsubsection*{Smart Transportation}

Random variables:

\begin{itemize}
\item $X$ = Vehicle Speed
\item $Y$ = Traffic Density
\end{itemize}

Question considered:

\begin{itemize}
\item Can vehicle speed be high when traffic density is high?
\end{itemize}

Observation:

\begin{itemize}
\item $X$ and $Y$ are \textbf{correlated random variables}.
\end{itemize}

Applications:

\begin{itemize}
\item Adaptive traffic signals
\item Congestion prediction
\item Autonomous vehicle routing
\end{itemize}

\subsubsection*{Wireless Network Performance}

Random variables:

\begin{itemize}
\item $X$ = Signal-to-Noise Ratio (SNR)
\item $Y$ = Data Throughput
\end{itemize}

Observation:

\begin{itemize}
\item Throughput depends on SNR.
\end{itemize}

Question considered:

\begin{itemize}
\item If SNR is high, will throughput always be high?
\end{itemize}

Reason:

\begin{itemize}
\item Network congestion affects throughput.
\end{itemize}

Applications:

\begin{itemize}
\item 5G scheduling
\item Adaptive modulation
\item Network planning
\end{itemize}

\subsubsection*{Weather Prediction}

Random variables:

\begin{itemize}
\item $X$ = Rainfall
\item $Y$ = Temperature
\end{itemize}

Application:

\begin{itemize}
\item Weather forecasting.
\end{itemize}

\subsubsection*{Drone Delivery}

Random variables:

\begin{itemize}
\item $X$ = Wind Speed
\item $Y$ = Delivery Time
\end{itemize}

Application:

\begin{itemize}
\item Logistics and UAV systems.
\end{itemize}

\subsubsection*{Dice Example}

Random variables:

\begin{itemize}
\item $X$ = outcome of first die
\item $Y$ = outcome of second die
\end{itemize}

Example conditions:

\begin{itemize}
\item $X + Y = 7$
\item $X > Y$
\end{itemize}

This requires \textbf{joint PMF analysis}.

\hrule
\vspace{0.5cm}

\section*{3. Core Results}

\subsection*{Discrete Case – Joint PMF}

\subsubsection*{Definition}

The \textbf{Joint Probability Mass Function (Joint PMF)} gives the probability that:

\[
X = x \quad \text{and} \quad Y = y
\]

\[
p_{X,Y}(x,y) = P(X = x, Y = y)
\]

This function represents the probability distribution for \textbf{two discrete random variables simultaneously}.

\subsection*{Joint PMF Table Representation}

The joint PMF for discrete variables can be represented using a \textbf{table}.

\begin{center}
\begin{tabular}{cccc}
\toprule
$Y \backslash X$ & $x_1$ & $x_2$ & $\dots$ \\
\midrule
$y_1$ & $p(x_1,y_1)$ & $p(x_2,y_1)$ & $\dots$ \\
$y_2$ & $p(x_1,y_2)$ & $p(x_2,y_2)$ & $\dots$ \\
\bottomrule
\end{tabular}
\end{center}

Each entry represents:

\[
P(X = x_i, Y = y_j)
\]

\subsection*{Continuous Case – Joint PDF}

For continuous random variables, probability is represented using the \textbf{Joint Probability Density Function (PDF)}.

\[
f_{X,Y}(x,y)
\]

The probability over a region is obtained by integration.

\hrule
\vspace{0.5cm}

\newpage
\section*{4. Analytical Development}

\subsection*{Geometric Interpretation (Continuous Case)}

For continuous variables:

\begin{itemize}
\item Probability corresponds to the \textbf{area under the joint density surface}.
\item The region in the $(x)-(y)$ plane determines the probability.
\end{itemize}

Thus,

\[
P((X,Y) \in R)
\]

is obtained by integrating the joint PDF over region $R$.

\hrule
\vspace{0.5cm}

\section*{5. Illustrative Applications}

\subsection*{Example: Rolling Two Dice}

Random variables:

\begin{itemize}
\item $X$ = result of die 1
\item $Y$ = result of die 2
\end{itemize}

Sample space size:

\[
6 \times 6 = 36
\]

Each ordered pair:

\[
(X,Y) = (1,1), (1,2), \dots, (6,6)
\]

Joint probabilities are represented in a \textbf{Joint PMF table}.

\hrule
\vspace{0.5cm}

\newpage
\section*{6. Joint Cumulative Distribution Function (Joint CDF)}

\subsection*{Definition}

The \textbf{Joint CDF} of two random variables $X$ and $Y$ is defined as:

\[
F_{X,Y}(x,y) = P(X \le x, Y \le y)
\]

This represents the probability that:

\begin{itemize}
\item $X$ is less than or equal to $x$
\item $Y$ is less than or equal to $y$
\end{itemize}

\subsection*{Continuous Case}

For continuous variables, the joint CDF is obtained from the joint PDF.

Conceptually:

\begin{itemize}
\item Integrate the joint density over the region
\end{itemize}

\[
(-\infty, x] \times (-\infty, y]
\]

\subsection*{Discrete Case}

For discrete random variables:

\[
F_{X,Y}(x,y)
=
\sum_{x_i \le x} \sum_{y_j \le y}
p_{X,Y}(x_i,y_j)
\]

The CDF is obtained by \textbf{summing the joint PMF values} up to $x$ and $y$.

\hrule
\vspace{0.5cm}

\section*{7. Properties of the Joint CDF}

The lecture introduces \textbf{properties of the Joint CDF}.

Key aspects:

\begin{itemize}
\item Defined for all real values of $x$ and $y$
\item Represents cumulative probability
\item Applicable to both discrete and continuous cases
\end{itemize}

\hrule
\vspace{0.5cm}

\section*{8. Structural Relationships}

The lecture follows the conceptual progression:

\begin{enumerate}
\item Motivation for joint random variables
\item Discrete case – Joint PMF
\item Joint PMF table representation
\item Dice example
\item Continuous case – Joint PDF
\item Geometric interpretation
\item Joint CDF definition
\item Continuous case of CDF
\item Discrete case of CDF
\item Properties of Joint CDF
\end{enumerate}

This structure builds from \textbf{basic motivation → mathematical formulation → examples → properties}.

\hrule
\vspace{0.5cm}

\newpage
\section*{9. Consolidated Summary}

\subsection*{Key Concepts}

\begin{itemize}
\item Joint random variables model \textbf{multiple random variables simultaneously}.
\item Relationships between variables may exist (correlation).
\end{itemize}

\subsection*{Important Functions}

\textbf{Joint PMF (Discrete Case)}

\[
p_{X,Y}(x,y) = P(X=x, Y=y)
\]

\textbf{Joint PDF (Continuous Case)}

\[
f_{X,Y}(x,y)
\]

\textbf{Joint CDF}

\[
F_{X,Y}(x,y) = P(X \le x, Y \le y)
\]

\subsection*{Example Contexts from the Lecture}

\begin{itemize}
\item Smart transportation
\item Wireless network performance
\item Weather prediction
\item Drone delivery
\item Rolling two dice
\end{itemize}

\end{document}